\section*{Chapitre \ref{ch:ll:lock-free-programming} --- Programmation sans verrou}

\begin{solution}{exo:ll:lock-free-programming:1}
Sur la mesure retenue, environ \qty{4.7}{\nano\second} seul et \qty{95}{\nano\second} avec quatre threads~; à quatre
threads, cela fait $4/\qty{95}{\nano\second} \approx 42$ millions d'incréments partagés par seconde, contre environ 730 millions avec
des compteurs par thread.
\end{solution}

\begin{solution}{exo:ll:lock-free-programming:2}
Avec un release sur le drapeau et un acquire sur sa lecture, non~: l'acquire qui voit le drapeau voit toutes les écritures faites avant le release,
y compris le prix. Avec un drapeau relaxed, oui~: rien n'ordonne l'écriture du prix avant celle du drapeau, ni pour le compilateur, ni
sur d'autres architectures, pour le processeur.
\end{solution}

\begin{solution}{exo:ll:lock-free-programming:3}
$2^{32}/10^8 \approx 43$ secondes. Un débordement ne produit une fausse correspondance que si un thread reste bloqué entre sa lecture et son CAS
pendant exactement un multiple de $2^{32}$ mises à jour et retrouve le même indice avec la même étiquette~: possible en théorie, extrêmement rare
en pratique, et écarté par une étiquette sur 64~bits avec un CAS sur 16~octets si besoin.
\end{solution}

\begin{solution}{exo:ll:lock-free-programming:4}
Sous \omterm{def:ll:lock-free-programming:ordering}{cohérence séquentielle}, $(0,1)$, $(1,0)$ et $(1,1)$~; sur x86 avec des ordonnancements plus faibles, $(0,0)$ est aussi possible, lorsque les deux écritures
se trouvent encore dans les tampons d'écriture de leurs cœurs au moment où les lectures se terminent. La fenêtre ne dure que quelques nanosecondes, alors que les départs des deux threads
diffèrent bien davantage à chaque essai~: elles se chevauchent donc rarement, soit quelques dizaines à quelques centaines de fois par million ici.
\end{solution}

\begin{solution}{exo:ll:lock-free-programming:5}
Un producteur réserve une position avec son \omterm{def:ll:lock-free-programming:cas}{compare-and-swap}, puis écrit la cellule et la publie. S'il est interrompu entre les deux, le numéro de séquence de la cellule
reste non publié~: le consommateur de cette position trouve la file «~vide~», et chaque position suivante attend derrière elle, si bien
qu'un seul thread bloqué arrête la progression des autres.
\end{solution}

\begin{solution}{exo:ll:lock-free-programming:6}
Le mutex est tenu 20~\% du temps ($100\,000 \times \qty{2}{\micro\second}$), donc environ un ordre sur cinq attend, une microseconde en moyenne.
Si le journaliseur est interrompu alors qu'il le tient, un ordre peut attendre tout le quantum, soit \qty{4}{\milli\second}.
\end{solution}

\begin{solution}{exo:ll:lock-free-programming:7} \sloppy
Sur vingt séries de deux millions d'opérations, aucun nœud n'a été perdu ni dupliqué (\texttt{measured\_aba.csv})~: la fenêtre vulnérable, entre la lecture du successeur
et le \omterm{def:ll:lock-free-programming:cas}{compare-and-swap}, ne dure que quelques nanosecondes, et l'ordonnanceur n'interrompt presque jamais un thread à l'intérieur.
Le code reste faux~; le rejeu déterministe de la \cref{fig:ll:lock-free-programming:aba} montre l'entrelacement.
Faire tourner du code \omterm{def:ll:lock-free-programming:progress}{sans verrou} prouve peu de chose~; l'argument doit venir du raisonnement, de la vérification de modèles ou d'outils
qui explorent les entrelacements.
\end{solution}

\begin{solution}{exo:ll:lock-free-programming:8}
La correction au sens du modèle C++ n'est pas une propriété d'un processeur donné~: x86 conserve par chance l'ordre des écritures, mais le compilateur peut
réordonner les écritures du snapshot après une écriture relaxed du drapeau, et une autre architecture ou une autre version du compilateur peut le faire. L'absence d'échec
lors d'un test n'est pas une preuve. Utiliser un release sur le drapeau et un acquire sur sa lecture, ou un verrou à séquence (chapitre~12).
\end{solution}

\begin{solution}{pb:ll:lock-free-programming:1}
\begin{enumerate}
\item $50\,000 \times \qty{3}{\micro\second} = 15~\%$ du temps.
\item 15~\% des ordres (les arrivées voient la moyenne temporelle), qui attendent en moyenne la moitié d'une tenue~: \qty{1.5}{\micro\second}.
\item $20\,000 \times 0.15 \times \qty{1.5}{\micro\second} = \qty{4.5}{\milli\second}$ par seconde.
\item Avec deux threads qui le prennent souvent, chaque acquisition coûte plusieurs fois son coût sans contention (la \omterm{def:ll:memory-hierarchy-and-caches:line}{ligne de cache} du verrou passe
d'un cœur à l'autre)~: environ \qty{100}{\nano\second} contre \qty{14}{\nano\second} sur la mesure retenue.
\item $50\,000 \times 0.0001 = 5$ par seconde.
\item $5 \times \qty{3}{\milli\second} = \qty{15}{\milli\second}$ par seconde~: 1.5~\% du temps.
\item $20\,000 \times 0.015 = 300$ ordres par seconde, chacun attendant au plus \qty{3}{\milli\second}, soit \qty{1.5}{\milli\second} en
moyenne.
\item $300 \times \qty{1.5}{\milli\second} + \qty{4.5}{\milli\second} \approx \qty{455}{\milli\second}$ par seconde~: les rares
interruptions dominent.
\item Il écrit l'ordre dans son propre emplacement du tampon circulaire (ring buffer) et le publie par une écriture release~: quelques dizaines de nanosecondes, sans
attendre le journaliseur.
\item La stratégie doit trancher~: abandonner l'enregistrement (et compter l'abandon) ou écraser, sans jamais bloquer~; le tampon circulaire est dimensionné de sorte qu'il
ne se remplisse pas.
\item \omterm{def:ll:lock-free-programming:progress}{Sans attente}~: chaque push se termine en un nombre borné de ses propres étapes.
\item L'héritage de priorité élève la priorité du journaliseur tant qu'il tient le verrou, mais l'échéance d'un thread de trading se mesure en microsecondes~;
réveiller et faire tourner le détenteur prend plus de temps que cela, et l'interruption par le noyau lui-même ne relève pas des priorités.
\item \textbf{Résultat nommé.} Avec le mutex, la stratégie est bloquée environ \qty{455}{\milli\second} par seconde en moyenne, presque tout
en attendant un journaliseur interrompu~; avec le tampon circulaire, jamais.
\item Ils ont activé l'héritage de priorité sur le mutex en déployant un correctif~: cela a suffi car leur échéance était
un cycle de bus de quelques millisecondes, et le détenteur pouvait finir vite une fois sa priorité relevée.
\item Dans la bibliothèque standard (allocateurs, journalisation, locales), dans les bibliothèques tierces, dans le chargeur dynamique, et dans tout
objet de configuration partagé.
\item Un moyen de savoir qu'aucun thread ne lit encore un nœud avant de le réutiliser~: des étiquettes sur les nœuds réutilisés, des \omterm{def:ll:lock-free-programming:reclaim}{pointeurs de danger} ou des époques sur
les nœuds libérés.
\item Deux producteurs et deux consommateurs font passer 300\,000 éléments et vérifient que chacun arrive une seule fois, en C++ et en Rust, et le test C++ tourne
sous le sanitiseur de threads. Ils ne peuvent pas prouver l'absence d'entrelacements rares.
\item Hors du chemin critique, pour un état touché rarement par des threads qui peuvent se permettre d'attendre (configuration, reporting), là où la simplicité
l'emporte.
\item Release sur l'écriture qui publie l'emplacement, et acquire sur la lecture qui le consomme.
\item Elle supprime l'attente sur les autres threads du chemin critique~: aucun blocage d'un thread ne peut l'arrêter.
\end{enumerate}
\end{solution}

\begin{solution}{iq:ll:lock-free-programming:1}
Relaxed~: atomicité seulement. Acquire--release~: une écriture release publie toutes les écritures antérieures au thread dont la lecture acquire lit
cette valeur. Séquentiellement cohérent~: en plus, un ordre total unique de toutes ces opérations, interdisant par exemple que les deux threads d'un test de
tampon d'écriture lisent zéro.
\iqlookfor{publication face à un ordre global, avec un exemple.}
\end{solution}

\begin{solution}{iq:ll:lock-free-programming:2}
Un \omterm{def:ll:lock-free-programming:cas}{compare-and-swap} réussit parce que la valeur est redevenue celle qu'il avait lue, bien qu'elle ait changé entre-temps et que l'état calculé à partir d'elle soit
périmé. On l'évite avec une étiquette incrémentée à chaque mise à jour (comparée au pointeur), ou en ne réutilisant jamais un nœud tant qu'un
lecteur peut le tenir (\omterm{def:ll:lock-free-programming:reclaim}{pointeurs de danger}, époques, ou nœuds jamais libérés).
\iqlookfor{le mécanisme et deux remèdes indépendants.}
\end{solution}

\begin{solution}{iq:ll:lock-free-programming:3}
Non. Sous forte contention sur un même emplacement, une boucle de CAS réessaie et peut être plus lente qu'un mutex~; les compteurs du chapitre montrent une boucle de CAS
à environ \qty{200}{\nano\second} par incrément avec quatre threads. Le caractère \omterm{def:ll:lock-free-programming:progress}{sans verrou} est une garantie sur les blocages, non sur la vitesse.
\iqlookfor{progression versus débit.}
\end{solution}

\begin{solution}{iq:ll:lock-free-programming:4}
Un thread de priorité élevée attend un verrou détenu par un thread de faible priorité, que des tâches de priorité moyenne empêchent de tourner. Sur un desk, le
thread critique attend un journaliseur ou un module de reporting que l'ordonnanceur a interrompu, pendant des millisecondes.
\iqlookfor{le mécanisme à trois threads et un exemple de trading.}
\end{solution}

\begin{solution}{iq:ll:lock-free-programming:5}
Des tests de résistance avec vérification des invariants (chaque élément une fois), des sanitiseurs, des tests qui perturbent l'ordonnancement, la vérification de modèles
de l'algorithme, et des rejeux déterministes d'entrelacements connus comme mauvais. Un test de résistance ne peut pas montrer qu'un entrelacement rare est
inoffensif, seulement qu'il ne s'est pas produit.
\iqlookfor{outils au-delà de l'exécution, et limites du test.}
\end{solution}

\begin{solution}{iq:ll:lock-free-programming:6}
Un compteur par thread, chacun sur sa propre \omterm{def:ll:memory-hierarchy-and-caches:line}{ligne de cache}, incrémenté par une opération relaxed (ou une simple écriture si seul son propriétaire écrit
et que le lecteur tolère une valeur légèrement périmée)~; le lecteur fait la somme. Aucune ligne n'est partagée entre écrivains.
\iqlookfor{sharding par thread et remplissage de \omterm{def:ll:memory-hierarchy-and-caches:line}{ligne de cache}.}
\end{solution}
